% Kompilacja z katalogu projektu:
% latexmk -pdf -interaction=nonstopmode -halt-on-error -outdir=tmp/instructions-update/build docs/instructions.tex
% Następnie: cp tmp/instructions-update/build/instructions.pdf docs/instructions.pdf
\documentclass[11pt,a4paper]{article}
\usepackage[T1]{fontenc}
\usepackage[utf8]{inputenc}
\usepackage[polish]{babel}
\usepackage{lmodern}
\usepackage[margin=22mm]{geometry}
\usepackage{array,booktabs,tabularx}
\usepackage{titlesec}
\usepackage[unicode,hidelinks]{hyperref}
\hypersetup{pdftitle={Uruchomienie Ponga},
  pdfauthor={Artur Zamorowski, Kacper Wojtowicz, Rafał Wybraniec},pdflang={pl-PL}}
\setlength{\parindent}{0pt}
\setlength{\parskip}{6pt}
\setlength{\emergencystretch}{1.5em}
\titleformat{\section}{\large\bfseries}{}{0pt}{}
\titlespacing*{\section}{0pt}{12pt}{5pt}
\pagestyle{empty}

\begin{document}
{\LARGE\bfseries Uruchomienie Ponga\par}

Wymagany jest Python 3.11 lub nowszy. Gotowe modele są dołączone.
Proszę rozpakować archiwum i otworzyć terminal w katalogu projektu,
w którym znajdują się \texttt{requirements.txt} i \texttt{config.yaml}.
Jeśli projekt został pobrany z GitHuba, proszę otworzyć katalog repozytorium.

\section*{1. Instalacja i uruchomienie}
Na macOS i Linux proszę wykonać:
\begin{verbatim}
python3 -m venv .venv
source .venv/bin/activate
python -m pip install -r requirements.txt
python src/main.py play
\end{verbatim}
Na Windows proszę wykonać:
\begin{verbatim}
py -3 -m venv .venv
.venv\Scripts\python.exe -m pip install -r requirements.txt
.venv\Scripts\python.exe src\main.py play
\end{verbatim}

\section*{2. Gra i sterowanie}
Domyślnie startuje pokaz Trackera przeciw agentowi, zgodnie z ustawieniem
\texttt{demo: true} w \texttt{config.yaml}. Aby zagrać samodzielnie,
proszę nacisnąć Esc, a następnie Enter. Mecz trwa do 7 punktów.

\renewcommand{\arraystretch}{1.15}
\begin{tabularx}{\linewidth}{@{}lX@{}}
\toprule
Klawisze & Czynność\\
\midrule
W/S lub strzałki & Proszę używać do poruszania lewą paletką.\\
Spacja & Proszę nacisnąć, aby zatrzymać lub wznowić grę.\\
D & Proszę nacisnąć, aby rozpocząć nowy mecz i przełączyć tryb pokazu.\\
R & Proszę nacisnąć, aby rozpocząć nowy mecz w bieżącym trybie.\\
H & Proszę nacisnąć, aby wyświetlić lub ukryć pomoc.\\
Esc & Proszę nacisnąć, aby wrócić do menu. Aby zamknąć grę z menu,
proszę nacisnąć Esc ponownie.\\
\bottomrule
\end{tabularx}

\section*{3. Pozostałe tryby}
W aktywnym środowisku wirtualnym proszę wykonać odpowiednie polecenie:
\begin{verbatim}
python src/main.py evaluate
python src/main.py matches
python src/main.py train
\end{verbatim}
\texttt{evaluate} ocenia model gry i baseline na tych samych serwisach.
\texttt{matches} rozgrywa pełne mecze z Trackerem bez okna gry.
\texttt{train} rozpoczyna nowy trening pojedynczego ustawienia
z sekcji \texttt{training}. Najlepsza tablica zastępuje model gry,
jeśli w ocenie końcowej wygrywa częściej, niż przegrywa.

Na Windows proszę zastąpić \texttt{python} przez
\texttt{.venv\textbackslash Scripts\textbackslash python.exe}.
Aby zmienić parametry lub wybrany model, proszę edytować \texttt{config.yaml}.
Ścieżki w konfiguracji odnoszą się do katalogu projektu.
\end{document}
